%\documentclass[a4paper]{article}
%\usepackage{geometry}
%\geometry{a4paper,scale=0.8}
\documentclass[10pt]{article}
\usepackage{ctex}
\usepackage{indentfirst}
\usepackage{longtable}
\usepackage{multirow}
\usepackage[a4paper, total={6in, 8in}]{geometry}
\usepackage{CJK}
\usepackage[fleqn]{amsmath}
\usepackage{parskip}
\usepackage{listings}
\usepackage{fancyhdr}

\setlength{\headheight}{15pt}
\pagestyle{fancy}

% 设置页眉
\fancyhead[L]{2026年秋季}
\fancyhead[C]{机器学习}
\fancyhead[R]{作业二}


\usepackage{graphicx}
\usepackage{float}
\usepackage{multicol}
\usepackage{amssymb}
\usepackage{booktabs}
\usepackage{xcolor}

% 定义Python代码风格
\definecolor{codegreen}{rgb}{0,0.6,0}
\definecolor{codegray}{rgb}{0.5,0.5,0.5}
\definecolor{codepurple}{rgb}{0.58,0,0.82}
\definecolor{backcolour}{rgb}{0.95,0.95,0.92}

\lstdefinestyle{mystyle}{
    backgroundcolor=\color{backcolour},   
    commentstyle=\color{codegreen},
    keywordstyle=\color{magenta},
    numberstyle=\tiny\color{codegray},
    stringstyle=\color{codepurple},
    basicstyle=\ttfamily\footnotesize,
    breakatwhitespace=false,         
    breaklines=true,     
    captionpos=b,        
    keepspaces=true,     
    numbers=left,        
    numbersep=5pt,       
    showspaces=false,    showstringspaces=false,
    showtabs=false,      
    tabsize=2
}

\lstset{style=mystyle}
\begin{document}

\textbf{\color{blue} \Large 姓名：闵振昌 \ \ \ 学号：241880334 \ \ \ \today}



\section*{一. (30 points) 线性回归与冗余属性}
考虑一元线性回归模型 $\hat y=wx+b$，最小化 $J(w,b)=\frac12\sum_{i=1}^{m}(y_i-wx_i-b)^2$。数据如下，记训练集均值为 $\bar x$ 与 $\bar y$。
\begin{center}
\begin{tabular}{ccccc}
\toprule
训练样本 & 1 & 2 & 3 & 4 \\
\midrule
$x$ & 0 & 1 & 2 & 3 \\
$y$ & 1 & 2 & 2 & 3 \\
\bottomrule
\end{tabular}
\qquad
\begin{tabular}{ccc}
\toprule
测试样本 & 1 & 2 \\
\midrule
$x$ & 4 & 5 \\
$y$ & 4 & 4 \\
\bottomrule
\end{tabular}
\end{center}
1. 对 $J$ 分别关于 $b$ 和 $w$ 求导，推导 $b^*=\bar y-w^*\bar x$ 与 $w^*$ 的表达式。假设 $\sum_i(x_i-\bar x)^2>0$。（10 points）

2. 求本题训练集上的 $w^*,b^*$，计算两个测试样本的预测值及测试集均方误差 $\mathrm{MSE}=\frac1n\sum_i(y_i-\hat y_i)^2$。（5 points）

3. 为每个样本增加属性 $x_2=2x$，模型变为 $\hat y=w_1x+w_2x_2+b$。说明 $(w_1,w_2)$ 是否唯一，并写出两组能得到与第2问相同预测的参数。测试样本也满足 $x_2=2x$。（5 points）

4. 补全以下代码中的两个函数，输出参数、测试预测值和 MSE。只能使用 NumPy 实现核心计算，不使用现成回归模型。（10 points）
\begin{lstlisting}[language=Python,caption=线性回归代码样例]
import numpy as np
X_train = np.array([[0.], [1.], [2.], [3.]])
y_train = np.array([1., 2., 2., 3.])
X_test = np.array([[4.], [5.]])
y_test = np.array([4., 4.])
def linear_regression(X_train, y_train):
    # X_train: (m, 1); y_train: (m,)
    pass  # Return scalar w and scalar b.
def MSE(y_true, y_pred):
    pass  # Return the mean squared error.
w, b = linear_regression(X_train, y_train)
y_pred = w * X_test[:, 0] + b
print(w, b, y_pred, MSE(y_test, y_pred))
\end{lstlisting}

\textbf{\large 解:}

\subsection*{1.}
分别对截距和斜率求偏导，令其等于零：
\begin{align*}
\frac{\partial J}{\partial b}
&=-\sum_{i=1}^{m}(y_i-wx_i-b)=0,\\
\frac{\partial J}{\partial w}
&=-\sum_{i=1}^{m}x_i(y_i-wx_i-b)=0.
\end{align*}
由第一个方程得
\[
\sum_{i=1}^{m}y_i-w\sum_{i=1}^{m}x_i-mb=0,
\qquad {b^*=\bar y-w^*\bar x}.
\]
将其代入第二个方程，得到
\[
\sum_{i=1}^{m}x_i\bigl[(y_i-\bar y)-w^*(x_i-\bar x)\bigr]=0.
\]
利用 $\sum_i(y_i-\bar y)=\sum_i(x_i-\bar x)=0$，整理得
\[
\sum_{i=1}^{m}(x_i-\bar x)(y_i-\bar y)
=w^*\sum_{i=1}^{m}(x_i-\bar x)^2.
\]
由题设 $\sum_i(x_i-\bar x)^2>0$，可得
\[
{w^*=\frac{\sum_{i=1}^{m}(x_i-\bar x)(y_i-\bar y)}
{\sum_{i=1}^{m}(x_i-\bar x)^2}},\qquad
{b^*=\bar y-w^*\bar x}.
\]

\subsection*{2.}
训练集均值为 $\bar x=1.5$、$\bar y=2$，且
\begin{align*}
\sum_i(x_i-\bar x)^2&=2.25+0.25+0.25+2.25=5,\\
\sum_i(x_i-\bar x)(y_i-\bar y)&=1.5+0+0+1.5=3.
\end{align*}
故 ${w^*=0.6,\ b^*=1.1}$，预测函数为 $\hat y=0.6x+1.1$。
两个测试样本的预测值分别为
\[
\hat y(4)=3.5,\qquad \hat y(5)=4.1.
\]
测试集均方误差为
\[
{\mathrm{MSE}=\frac{(4-3.5)^2+(4-4.1)^2}{2}=0.13}.
\]

\subsection*{3.}
由于 $x_2=2x$，新模型可写为
\[
\hat y=(w_1+2w_2)x+b.
\]
只要 $w_1+2w_2=0.6$ 且 $b=1.1$，就与第2问得到相同的预测，
因此 $(w_1,w_2)$ 不唯一。两组参数为
\[
{(w_1,w_2,b)=(0.6,0,1.1)},\qquad
{(w_1,w_2,b)=(0,0.3,1.1)}.
\]
测试样本同样满足 $x_2=2x$，所以这两组参数的预测值均为 $3.5$、$4.1$。

\subsection*{4.}
\begin{lstlisting}[language=Python,caption=一元线性回归与均方误差的完整实现]
import numpy as np

X_train = np.array([[0.], [1.], [2.], [3.]])
y_train = np.array([1., 2., 2., 3.])
X_test = np.array([[4.], [5.]])
y_test = np.array([4., 4.])

def linear_regression(X_train, y_train):
    x = np.asarray(X_train, dtype=float)[:, 0]
    y = np.asarray(y_train, dtype=float)
    x_mean = np.mean(x)
    y_mean = np.mean(y)
    dx = x - x_mean
    denominator = np.sum(dx ** 2)
    w = np.sum(dx * (y - y_mean)) / denominator
    b = y_mean - w * x_mean
    return float(w), float(b)

def MSE(y_true, y_pred):
    y_true = np.asarray(y_true, dtype=float)
    y_pred = np.asarray(y_pred, dtype=float)
    return float(np.mean((y_true - y_pred) ** 2))

w, b = linear_regression(X_train, y_train)
y_pred = w * X_test[:, 0] + b
print(w, b, y_pred, MSE(y_test, y_pred))
\end{lstlisting}
输出为：
\begin{lstlisting}
0.6 1.1 [3.5 4.1] 0.12999999999999998
\end{lstlisting}
即 $w=0.6$、$b=1.1$，预测值为 $3.5$、$4.1$，$\mathrm{MSE}=0.13$。


\newpage
\section*{二. (25 points) 对数几率回归与类别不平衡}
某二分类任务规定少数类为正类 $y=1$，多数类为负类 $y=0$。对数几率回归输出正类概率
\[
p(\mathbf{x})=\frac{1}{1+\exp[-(\textbf{w}^{\top}\mathbf{x}+b)]}.
\]
给定已训练模型在10个评价样本上的输出如下。此模型的训练集共有100个样本，其中正类20个、负类80个；下表不是训练集。
\begin{center}
\begin{tabular}{c*{10}{c}}
\toprule
编号 & 1 & 2 & 3 & 4 & 5 & 6 & 7 & 8 & 9 & 10 \\
\midrule
$y$ & 0 & 0 & 0 & 0 & 0 & 0 & 0 & 0 & 1 & 1 \\
$p$ & .05 & .10 & .15 & .20 & .25 & .30 & .35 & .55 & .40 & .80 \\
\bottomrule
\end{tabular}
\end{center}
\textbf{规定：}使用阈值 $t$ 时，$p\geq t$ 判为正类，否则判为负类。所有问题均保持模型参数不变。

1. 当阈值为 $0.5$ 时，将预测规则等价改写为关于 $\mathbf{w}^{\top}\mathbf{x}+b$ 的不等式（消去概率 $p$），并说明该模型决策边界的形状。（5分）


2. 分别在阈值 $0.5$ 与 $0.3$ 下，计算 $TP,FP,TN,FN$ 及准确率、查准率、查全率。说明降低阈值对本题结果的影响。（8 points）

3. 按教材第3.6节的再缩放思想，本题规定当
\[
\frac{p}{1-p}\geq\frac{m^+}{m^-}
\]
时判为正类，其中 $m^+,m^-$ 分别为训练集正、负例数。推导等价的概率阈值，求其在本题的数值。简述这种固定模型调整阈值的方法与欠采样后重新训练的区别，不要求二者效果相同。（7 points）

4. 补全以下函数，返回第2问所需的计数和指标，并用两个阈值核对手算结果。（5 points）
\begin{lstlisting}[language=Python,caption=分类评价代码框架]
import numpy as np
y = np.array([0, 0, 0, 0, 0, 0, 0, 0, 1, 1])
p = np.array([.05, .10, .15, .20, .25,
              .30, .35, .55, .40, .80])
def evaluate(y, p, threshold):
    # Return: TP, FP, TN, FN, accuracy, precision, recall.
    # If a metric denominator is zero, return 0.
    pass
for t in [0.5, 0.3]:
    print(t, evaluate(y, p, t))
\end{lstlisting}
\noindent 提示：准确率 $=\frac{TP+TN}{TP+FP+TN+FN}$，查准率 $=\frac{TP}{TP+FP}$，查全率 $=\frac{TP}{TP+FN}$。

\textbf{\large 解:}

\subsection*{1.}
记 $z=\mathbf w^{\top}\mathbf x+b$。则
\[
p\geq\frac12
\iff \frac{1}{1+e^{-z}}\geq\frac12
\iff e^{-z}\leq1
\iff z\geq0.
\]
所以 $\mathbf w^{\top}\mathbf x+b\geq0$ 时预测正类，否则预测负类。
当 $\mathbf w\neq\mathbf0$ 时，决策边界是
$\mathbf w^{\top}\mathbf x+b=0$，在二维空间为直线，在高维空间为超平面。

\subsection*{2.}
当 $t=0.5$ 时，预测为正类的是样本8、10。
样本10是真正例，样本8是假正例，样本9是假负例，其余7个样本是真负例。
因此
\[
TP=1,\quad FP=1,\quad TN=7,\quad FN=1,
\]
\[
\mathrm{Accuracy}=\frac{1+7}{10}=0.8,\qquad
\mathrm{Precision}=\frac{1}{2}=0.5,\qquad
\mathrm{Recall}=\frac{1}{2}=0.5.
\]
当 $t=0.3$ 时，样本6、7、8、9、10预测为正类；其中概率恰为
$0.30$ 的样本6也包括在内。故
\[
TP=2,\quad FP=3,\quad TN=5,\quad FN=0,
\]
\[
\mathrm{Accuracy}=\frac{2+5}{10}=0.7,\qquad
\mathrm{Precision}=\frac{2}{5}=0.4,\qquad
\mathrm{Recall}=\frac{2}{2}=1.
\]
本题降低阈值后，多识别了1个正例，同时增加了2个假正例。
查全率由 $0.5$ 提升至 $1$，查准率由 $0.5$ 降至 $0.4$，准确率由 $0.8$ 降至 $0.7$。

\subsection*{3.}
对于 $0<p<1$，有
\begin{align*}
\frac{p}{1-p}\geq\frac{m^+}{m^-}
&\iff m^-p\geq m^+(1-p)\\
&\iff (m^++m^-)p\geq m^+\\
&\iff p\geq\frac{m^+}{m^++m^-}.
\end{align*}
代入训练集正、负例数 $m^+=20$、$m^-=80$，得到等价概率阈值
\[
{t=\frac{20}{20+80}=0.2}.
\]
固定模型调整阈值只改变类别判定规则，模型参数和输出概率保持不变。
欠采样后重新训练会减少多数类训练样本，并重新估计模型参数，输出概率也可能改变。
因此两种方法的结果不一定相同。

\subsection*{4.}
\begin{lstlisting}[language=Python,caption=二分类计数与评价指标的完整实现]
import numpy as np

y = np.array([0, 0, 0, 0, 0, 0, 0, 0, 1, 1])
p = np.array([.05, .10, .15, .20, .25,
              .30, .35, .55, .40, .80])

def evaluate(y, p, threshold):
    y = np.asarray(y)
    p = np.asarray(p, dtype=float)
    pred_positive = p >= threshold
    actual_positive = y == 1
    actual_negative = y == 0
    TP = int(np.sum(pred_positive & actual_positive))
    FP = int(np.sum(pred_positive & actual_negative))
    TN = int(np.sum(~pred_positive & actual_negative))
    FN = int(np.sum(~pred_positive & actual_positive))

    total = TP + FP + TN + FN
    accuracy = (TP + TN) / total if total else 0.0
    precision = TP / (TP + FP) if TP + FP else 0.0
    recall = TP / (TP + FN) if TP + FN else 0.0
    return TP, FP, TN, FN, accuracy, precision, recall

for t in [0.5, 0.3]:
    print(t, evaluate(y, p, t))
\end{lstlisting}
输出为：
\begin{lstlisting}
0.5 (1, 1, 7, 1, 0.8, 0.5, 0.5)
0.3 (2, 3, 5, 0, 0.7, 0.4, 1.0)
\end{lstlisting}
与第2问的手算结果一致。


\newpage
\section*{三. (20 points) 连续属性划分与缺失值}
二分类数据集 $D$ 含8个样本，连续属性 $x$ 与类别 $y$ 如下。每个样本的初始权重均为1；$\tilde D$ 表示属性 $x$ 不缺失的样本集合。对数以2为底，结果保留三位小数。
\begin{center}
\begin{tabular}{c*{8}{c}}
\toprule
编号 & 1 & 2 & 3 & 4 & 5 & 6 & 7 & 8 \\
\midrule
$x$ & 1 & 2 & 3 & 4 & 5 & 6 & 缺失 & 缺失 \\
$y$ & 0 & 0 & 1 & 1 & 0 & 1 & 0 & 1 \\
\bottomrule
\end{tabular}
\end{center}
按 $x\leq t$ 和 $x>t$ 将非缺失样本划分为左、右两个分支。

1. 列出由相邻不同属性值的中点构成的全部候选阈值，计算非缺失样本集合的熵 $H(\tilde D)$。（5 points）

2. 为减少计算，本问只比较 $t=2.5$ 与 $t=4.5$。分别计算它们在 $\tilde D$ 上的信息增益，并从这两个阈值中选出更优者。（5 points）

3. 考虑缺失值时，按教材方法使用
\[
\mathrm{Gain}(D,x,t)=\rho\,\mathrm{Gain}(\tilde D,x,t),
\qquad \rho=\frac{|\tilde D|}{|D|}.
\]
计算第2问两个阈值修正后的信息增益，判断选择是否改变。（5 points）

4. 使用第2问选出的阈值。按非缺失样本在各分支的权重占比，将每个缺失样本以相应权重同时下传到左右分支。写出样本7和8各自在左右分支的权重，并计算分配后左右分支的正、负类权重总和。（5 points）

\noindent 提示：$H(S)=-\sum_{c\in\{0,1\}}p_c\log_2p_c$，规定 $0\log_2 0=0$。第2问只在非缺失样本上计算
\[
\mathrm{Gain}(\tilde D,x,t)=H(\tilde D)
-\frac{|\tilde D_L|}{|\tilde D|}H(\tilde D_L)
-\frac{|\tilde D_R|}{|\tilde D|}H(\tilde D_R).
\]

\textbf{\large 解:}

\subsection*{1.}
非缺失属性值依次为 $1,2,3,4,5,6$，因此全部候选阈值为
\[
{\{1.5,\ 2.5,\ 3.5,\ 4.5,\ 5.5\}}.
\]
$\tilde D$ 共6个样本，正、负类各3个，故
\[
{H(\tilde D)=-\frac12\log_2\frac12-\frac12\log_2\frac12=1.000}.
\]

\subsection*{2.}
当 $t=2.5$ 时，左分支为样本1、2，类别均为0，故
$H(\tilde D_L)=0$；右分支为样本3、4、5、6，含3个正例和1个负例，故
\[
H(\tilde D_R)=-\frac34\log_2\frac34-\frac14\log_2\frac14
\approx0.811.
\]
于是
\[
\mathrm{Gain}(\tilde D,x,2.5)
=1-\frac26\times0-\frac46H(\tilde D_R)
\approx{0.459}.
\]
当 $t=4.5$ 时，左分支为样本1、2、3、4，正、负例各2个；
右分支为样本5、6，正、负例各1个。因此两个分支的熵均为 $1.000$，
\[
\mathrm{Gain}(\tilde D,x,4.5)
=1-\frac46\times1-\frac26\times1={0.000}.
\]
只在本题要求比较的两个阈值中，${t=2.5}$ 更优。

\subsection*{3.}
非缺失样本所占比例为
\[
\rho=\frac{|\tilde D|}{|D|}=\frac68=0.750.
\]
修正后的信息增益为
\begin{align*}
\mathrm{Gain}(D,x,2.5)
&=\frac68\,\mathrm{Gain}(\tilde D,x,2.5)\approx{0.344},\\
\mathrm{Gain}(D,x,4.5)
&=\frac68\times0={0.000}.
\end{align*}
同一属性的两个候选阈值均乘以相同的正数 $\rho$，
所以二者优劣次序不变，仍选择 $t=2.5$。

\subsection*{4.}
在 $t=2.5$ 时，非缺失样本的左、右分支权重分别为2、4，
所以缺失样本分配到左右分支的比例分别为 $1/3$、$2/3$。
样本7和8的初始权重均为1，因此
\begin{center}
\begin{tabular}{c|c|cc}
\toprule
样本 & 类别 & 左分支权重 & 右分支权重\\
\midrule
7 & 0 & $1/3\approx0.333$ & $2/3\approx0.667$\\
8 & 1 & $1/3\approx0.333$ & $2/3\approx0.667$\\
\bottomrule
\end{tabular}
\end{center}
加入这些样本后，各分支的类别权重为
\begin{align*}
W_L^+&=0+\frac13=\frac13\approx0.333,&
W_L^-&=2+\frac13=\frac73\approx2.333,\\
W_R^+&=3+\frac23=\frac{11}{3}\approx3.667,&
W_R^-&=1+\frac23=\frac53\approx1.667.
\end{align*}


\newpage
\section*{四. (25 points) 决策树的预剪枝与后剪枝}
某数据集有两个二值属性 $A,B\in\{0,1\}$。训练集与验证集的样本统计如下，每行表示相同属性取值下各类别的样本数量。
\begin{center}
\begin{tabular}{ccrrrr}
\toprule
\multicolumn{2}{c}{属性} & \multicolumn{2}{c}{训练集} & \multicolumn{2}{c}{验证集} \\
$A$ & $B$ & $y=0$ & $y=1$ & $y=0$ & $y=1$ \\
\midrule
0 & 0 & 7 & 0 & 3 & 0 \\
0 & 1 & 0 & 3 & 0 & 2 \\
1 & 0 & 0 & 4 & 0 & 2 \\
1 & 1 & 2 & 0 & 3 & 0 \\
\bottomrule
\end{tabular}
\end{center}
本题\textbf{固定划分顺序}：先在根节点按 $A$ 划分，再在两个子节点分别按 $B$ 划分，不要求比较划分指标。叶节点的预测类别始终由到达该节点的\textbf{训练样本多数类}决定；若数量相同则预测0。

1. 写出根节点不划分时的预测类别、按 $A$ 划分后两个叶节点的预测类别，以及完整两层树四个叶节点的预测类别。（5 points）

2. 预剪枝规则为：只有划分后\textbf{整棵树的验证准确率严格提高}，才允许该划分。计算根节点划分前、后的验证准确率，确定预剪枝结果。（7 points）

3. 从完整两层树开始进行后剪枝，顺序为：先检查 $A=0$ 子树，再检查 $A=1$ 子树，最后检查根节点。每次把待检查子树替换为一个多数类叶节点；若整棵树的验证准确率不下降，就保留该剪枝，否则恢复子树。计算各次候选剪枝后的验证准确率，画出或用清晰规则表示最终树。（8 points）

4. 结合本题解释为什么预剪枝和后剪枝可能得到不同结果，并说明验证集与测试集在剪枝实验中的不同用途。（5 points）

\noindent 提示：训练集共16个样本，验证集共10个样本。后剪枝被拒绝后，下一步仍以恢复后的树为基础。本题只需手算与分析，不要求另选数据集或实现剪枝算法。

\textbf{\large 解:}

\subsection*{1.}
根节点训练样本中，负例数为 $7+2=9$，正例数为 $3+4=7$，
因此不划分时预测类别为 ${0}$。
根节点按 $A$ 划分后：
\begin{itemize}
\item $A=0$：负例7个、正例3个，叶节点预测0；
\item $A=1$：负例2个、正例4个，叶节点预测1。
\end{itemize}
继续分别按 $B$ 划分后，四个叶节点的预测为
\begin{center}
\begin{tabular}{cc|rr|c}
\toprule
$A$ & $B$ & 训练负例数 & 训练正例数 & 预测类别\\
\midrule
0 & 0 & 7 & 0 & 0\\
0 & 1 & 0 & 3 & 1\\
1 & 0 & 0 & 4 & 1\\
1 & 1 & 2 & 0 & 0\\
\bottomrule
\end{tabular}
\end{center}

\subsection*{2.}
根节点不划分时，所有验证样本均预测为0，正确数为 $3+3=6$，所以
\[
\mathrm{Acc}_{\text{不划分}}=\frac6{10}=60\%.
\]
根节点按 $A$ 划分后，$A=0$ 的叶节点预测0，在对应验证样本中正确3个；
$A=1$ 的叶节点预测1，在对应验证样本中正确2个。因此
\[
\mathrm{Acc}_{\text{按 }A\text{ 划分}}=\frac{3+2}{10}=50\%.
\]
准确率从 $60\%$ 降到 $50\%$，不满足严格提高的条件，故拒绝根节点划分。
最终预剪枝树仅含一个预测0的根叶节点，验证准确率为 ${60\%}$，
也不会再尝试其下的 $B$ 划分。

\subsection*{3.}
完整两层树的四个叶节点都能正确分类对应的验证样本，初始验证准确率为
\[
\mathrm{Acc}_{\text{完整树}}=\frac{3+2+2+3}{10}=100\%.
\]
按题目顺序逐一检查：

\textbf{检查 $A=0$ 子树：}
将其替换为由训练多数类决定的类别0叶节点。
此时 $A=0$ 的5个验证样本中正确3个，未剪的 $A=1$ 子树仍正确5个，故
\[
\mathrm{Acc}_{\text{候选剪枝 }A=0}=\frac{3+5}{10}=80\%.
\]
低于当前的 $100\%$，拒绝此次剪枝并恢复完整树。

\textbf{检查 $A=1$ 子树：}
从恢复后的完整树出发，将该子树替换为预测1的叶节点。
$A=0$ 子树正确5个，$A=1$ 的叶节点正确2个，故
\[
\mathrm{Acc}_{\text{候选剪枝 }A=1}=\frac{5+2}{10}=70\%.
\]
仍低于 $100\%$，拒绝此次剪枝并恢复完整树。

\textbf{检查根节点：}
将整棵树替换为预测0的叶节点，得到
\[
\mathrm{Acc}_{\text{候选剪枝根节点}}=\frac6{10}=60\%.
\]
低于 $100\%$，故也拒绝此次剪枝。
最终保留完整两层树，验证准确率为 ${100\%}$。
按先 $A$、再 $B$ 的顺序，其明确预测规则为
\[
{\hat y(A,B)=
\begin{cases}
0,& A=0,\ B=0,\\
1,& A=0,\ B=1,\\
1,& A=1,\ B=0,\\
0,& A=1,\ B=1.
\end{cases}}
\]

\subsection*{4.}
本题按 $A$ 划分后，验证准确率由 $60\%$ 降为 $50\%$，预剪枝因此停止划分，
得到预测0的单叶树，无法继续考察下一层按 $B$ 划分带来的收益。
后剪枝先生成完整两层树，其验证准确率为 $100\%$，三次候选剪枝均使准确率下降，
因此保留完整树。这说明当前一步划分效果不好，并不代表后续划分不能改善结果，
所以预剪枝和后剪枝可能得到不同结果。

验证集用于比较剪枝前后的准确率、决定是否剪枝；测试集用于在最终树确定后
评价其泛化性能，不参与剪枝决策。


\end{document}
